% Kodowanie i czcionka
\usepackage[utf8]{inputenc}         % Kodowanie UTF-8
\usepackage[T1]{fontenc}            % Lepsze dzielenie wyrazów i znaki diakrytyczne
\usepackage{helvet}                 % Czcionka Helvetica
\renewcommand{\familydefault}{\sfdefault} % Domyślna czcionka bezszeryfowa

% Matematyka i symbole
\usepackage{amsmath,amsfonts,amssymb,amsthm} % Rozszerzenia matematyczne


% Marginesy
%\usepackage[margin=2.5cm]{geometry} % Ustawienie marginesów
\usepackage[top=2.5cm, bottom=2.5cm, left=3.5cm, right=1.5cm]{geometry} % lewy i prawwy inaczej dla marginesow 


% Ulepszenia typografii
\usepackage{microtype}
\usepackage{xcolor}
\usepackage{array}
\usepackage{enumitem} % opisy do wzorow
% ODSTEPY ENUMERATE 1. 2. 3. 
\setlist[enumerate]{itemsep=0pt, parsep=0pt, topsep=0pt}

% Odnośniki i linki
\usepackage{url}
\usepackage[hidelinks]{hyperref}   % Linki bez kolorów/ramki
\urlstyle{same} % linki taka sama czcionka



% --- FORMATOWANIE NAGŁÓWKÓW (Musi być dodane) ---
\usepackage{titlesec}

% Formatowanie ROZDZIAŁÓW (Nagłówek 1)
% Wymogi: 16 pkt, pogrubiona, odstępy 18pkt przed i po
\titleformat{\chapter}[block]
  {\bfseries\fontsize{16}{16}\selectfont} % Styl: Pogrubienie, rozmiar 16pt
  {\chaptername\ \thechapter.} % Numeracja: "Rozdział X."
  {0.5em} % Odstęp między "Rozdział X." a tytułem
  {}

% Odstępy dla rozdziału: {lewy}{góra}{dół}
\titlespacing*{\chapter}{0pt}{-20pt}{18pt}

% Formatowanie ROZDZIAŁÓW nienumerowanych (Spisy, Literatura, Wstęp)
\titleformat{name=\chapter,numberless}[block]
  {\bfseries\fontsize{16}{16}\selectfont}
  {}
  {0pt}
  {}
\titlespacing*{name=\chapter,numberless}{0pt}{-20pt}{18pt}

% Formatowanie PODROZDZIAŁÓW (Nagłówek 2 - np. 1.1)
% Wymogi: 14 pkt, pogrubiona, odstępy 12pkt przed i po
\titleformat{\section}[block]
  {\bfseries\fontsize{14}{14}\selectfont}
  {\thesection.}
  {0.5em}
  {}
\titlespacing*{\section}{0pt}{12pt}{12pt}

% Formatowanie PODPUNKTÓW (Nagłówek 3 - np. 1.1.1)
% Wymogi: 12 pkt, pogrubiona, odstępy 6pkt przed i po
\titleformat{\subsection}[block]
  {\bfseries\fontsize{12}{12}\selectfont}
  {\thesubsection.}
  {0.5em}
  {}
\titlespacing*{\subsection}{0pt}{6pt}{6pt}



% Wcięcia i bibliografia ************************************************************************************
\usepackage{indentfirst}           % Wcięcie w pierwszym akapicie każdego rozdziału
\setlength{\parindent}{1cm}

% --- BIBLIOGRAFIA I CYTOWANIA (Poprawiona) ---
\usepackage[polish]{babel}
\usepackage{csquotes}
\MakeOuterQuote{"}

% Konfiguracja BibLaTeX - styl authoryear (APA-like)
\usepackage[
    style=authoryear,    % Styl: Autor (Rok)
    backend=biber,       % Silnik przetwarzania
    maxcitenames=2,      % W tekście: max 2 nazwiska, potem "i in." (zgodnie z APA)
    maxbibnames=99,      % W bibliografii: wymień wszystkich autorów
    uniquelist=false,    % Nie dodawaj inicjałów, jeśli nazwiska się powtarzają
    isbn=false,          % WYŁĄCZ ISBN (zgodnie z szablonem)
    doi=false,           % WYŁĄCZ DOI (zgodnie z szablonem)
    url=false,           % WYŁĄCZ URL dla książek/artykułów
    eprint=false
]{biblatex}

% Włączamy URL tylko dla typu @online (netografia)
\AtEveryBibitem{%
  \ifentrytype{online}{}{\clearfield{url}}%
  \ifentrytype{online}{}{\clearfield{urldate}}%
}

% --- FORMATOWANIE WYGLĄDU bibliografi ---

% 1. Wszystkie tytuły kursywą (książki, artykuły, strony www)
\DeclareFieldFormat{title}{\mkbibemph{#1}}          
\DeclareFieldFormat[article]{title}{\mkbibemph{#1}} 
\DeclareFieldFormat[book]{title}{\mkbibemph{#1}}    
\DeclareFieldFormat[online]{title}{\mkbibemph{#1}}    

% 2. USUWANIE "W:" przed nazwą czasopisma
% Standardowo biblatex daje "In:" lub "W:". Ta komenda to usuwa.
\renewbibmacro{in:}{}

% 3. Kolejność: Wydawnictwo, Miasto (Dla książek)
% Standard to Miasto: Wydawnictwo. UE Katowice chce odwrotnie.
\renewbibmacro*{publisher+location+date}{%
  \printlist{publisher}%
  \setunit*{\addcomma\space}%
  \printlist{location}%
  \setunit*{\addcomma\space}%
  \usebibmacro{date}%
  \newunit}

% 4. Separatory
\renewcommand*{\nameyeardelim}{\addcomma\space} % Przecinek po nazwisku w cytowaniu
\renewcommand*{\postnotedelim}{\addcomma\space} % Przecinek przed numerem strony

% 5. Data dostępu dla netografii: [dostęp DD.MM.RRRR]
\DefineBibliographyStrings{polish}{
  urlseen = {dostęp},
}
\DeclareFieldFormat{urldate}{\mkbibbrackets{\bibstring{urlseen}\space#1}}
\DeclareFieldFormat{url}{\url{#1}}

% 6. LISTA NUMEROWANA (1. 2. 3.) do bibliografi 
\newcounter{mybibcounter}
\defbibenvironment{bibliography}
  {\list
     {\stepcounter{mybibcounter}\themybibcounter.} 
     {\setlength{\leftmargin}{2.5em}%              
      \setlength{\labelwidth}{2em}%                
      \setlength{\labelsep}{0.5em}%                
      \setlength{\itemindent}{0pt}%
      \setlength{\itemsep}{\bibitemsep}%
      \setlength{\parsep}{\bibparsep}}}
  {\endlist}
  {\item}





% Kod źródłowy (np. programy)
\usepackage{listings}

% Nagłówki i stopki
\usepackage{fancyhdr}
\usepackage{etoolbox}

% Grafika
\usepackage{graphicx}
\graphicspath{ {images/} }         % Folder z grafikami

% Opisy pod rysunkami
\usepackage{caption}
\usepackage{chngcntr}

% Globalne numerowanie rysunków i tabel
\counterwithout{figure}{chapter}
\counterwithout{table}{chapter}

%definicja podpisu obrazka
\DeclareCaptionFont{elevenpt}{\fontsize{11}{13.6}\selectfont} % Definicja 11pkt
\captionsetup{
  font=elevenpt, % Użycie 11pkt
  labelfont=bf,
  labelsep=period,
  justification=raggedright, 
  singlelinecheck=false
}


\usepackage{float}

\newcommand{\MyFigure}[5][0.7\textwidth]{%
  \begin{figure}[H] % H dla tego miejsca ht dla ustawienia roznie 
    \centering
    \includegraphics[width=#1]{#2}
    \caption{#3}
    \label{#4}
    \vspace{-3pt}
    {\raggedright\fontsize{10pt}{12pt}\selectfont\textit{Źródło: #5}\par}
  \end{figure}
}

% Niestandardowe środowisko do tabel z podpisem i źródłem ---- TABELE
\captionsetup[table]{
    format=plain,
    justification=raggedright,
   labelfont=bf,
   textfont=normalfont,
   labelsep=period,
   name=Tabela,
   singlelinecheck=false,
   skip=2pt
}

 \newcommand{\MyTableSource}[1]{%
   \vspace{-8pt}
   \begin{flushleft}
     {\linespread{0.9}\fontsize{10}{12pt}\selectfont\parbox{0.95\textwidth}{\raggedright\textit{Źródło: #1}}}
   \end{flushleft}
   \vspace{-6pt}
 }

% Nadpisanie stylu rozdziału – aby używał fancyhdr
\patchcmd{\chapter}{\thispagestyle{plain}}{\thispagestyle{fancy}}{}{}

% Stylowanie kodu źródłowego
\definecolor{codegreen}{rgb}{0,0.6,0}
\definecolor{codegray}{rgb}{0.5,0.5,0.5}
\definecolor{codepurple}{rgb}{0.58,0,0.82}
\definecolor{backcolour}{rgb}{0.98,0.98,0.95} % Nieco jaśniejsze tło

\lstdefinestyle{mystyle}{
    backgroundcolor=\color{backcolour},
    commentstyle=\color{codegreen},
    keywordstyle=\color{blue}\bfseries, % Słowa kluczowe na niebiesko i pogrubione
    numberstyle=\tiny\color{codegray},
    stringstyle=\color{codepurple},
    basicstyle=\ttfamily\footnotesize,
    breakatwhitespace=false,
    breaklines=true,
    captionpos=b,
    keepspaces=true,
    numbers=left,
    numbersep=10pt,                  % Większy odstęp numerów od kodu
    showspaces=false,
    showstringspaces=false,
    showtabs=false,
    tabsize=4,                       % Standard dla Pythona to 4 spacje
    frame=single,                    % Ramka wokół kodu
    rulecolor=\color{black!10},      % Kolor ramki
    extendedchars=true,              % Obsługa znaków specjalnych
    literate={ą}{{\k{a}}}1 {ć}{{\'c}}1 {ę}{{\k{e}}}1 {ł}{{\l{}}}1 {ń}{{\'n}}1 {ó}{{\'o}}1 {ś}{{\'s}}1 {ź}{{\'z}}1 {ż}{{\.z}}1 {Ą}{{\k{A}}}1 {Ć}{{\'C}}1 {Ę}{{\k{E}}}1 {Ł}{{\L{}}}1 {Ń}{{\'N}}1 {Ó}{{\'O}}1 {Ś}{{\'S}}1 {Ź}{{\'Z}}1 {Ż}{{\.Z}}1 % Polskie znaki w komentarzach
}

% Definicja konkretnie pod Pythona
\lstdefinestyle{pythonstyle}{
    style=mystyle,
    language=Python,
    morekeywords={self, cls, @property, @staticmethod, yield}
}

\lstset{style=pythonstyle} % Ustawienie jako domyślny

\lstdefinelanguage{javaScript}{
  keywords={typeof, new, true, false, catch, function, return, null, catch, switch, var, if, in, while, do, else, case, break},
  keywordstyle=\color{blue}\bfseries,
  ndkeywords={class, export, boolean, throw, implements, import, this},
  ndkeywordstyle=\color{darkgray}\bfseries,
  identifierstyle=\color{black},
  sensitive=false,
  comment=[l]{//},
  morecomment=[s]{/*}{*/},
  commentstyle=\color{purple}\ttfamily,
  stringstyle=\color{red}\ttfamily,
  morestring=[b]',
  morestring=[b]"
}

\lstdefinelanguage{css}{
  morekeywords={accelerator,azimuth,background,background-attachment,
    background-color,background-image,background-position,
    background-position-x,background-position-y,background-repeat,
    behavior,border,border-bottom,border-bottom-color,
    border-bottom-style,border-bottom-width,border-collapse,
    border-color,border-left,border-left-color,border-left-style,
    border-left-width,border-right,border-right-color,
    border-right-style,border-right-width,border-spacing,
    border-style,border-top,border-top-color,border-top-style,
    border-top-width,border-width,bottom,caption-side,clear,
    clip,color,content,counter-increment,counter-reset,cue,
    cue-after,cue-before,cursor,direction,display,elevation,
    empty-cells,filter,float,font,font-family,font-size,
    font-size-adjust,font-stretch,font-style,font-variant,
    font-weight,height,ime-mode,include-source,
    layer-background-color,layer-background-image,layout-flow,
    layout-grid,layout-grid-char,layout-grid-char-spacing,
    layout-grid-line,layout-grid-mode,layout-grid-type,left,
    letter-spacing,line-break,line-height,list-style,
    list-style-image,list-style-position,list-style-type,margin,
    margin-bottom,margin-left,margin-right,margin-top,
    marker-offset,marks,max-height,max-width,min-height,
    min-width,-moz-binding,-moz-border-radius,
    -moz-border-radius-topleft,-moz-border-radius-topright,
    -moz-border-radius-bottomright,-moz-border-radius-bottomleft,
    -moz-border-top-colors,-moz-border-right-colors,
    -moz-border-bottom-colors,-moz-border-left-colors,-moz-opacity,
    -moz-outline,-moz-outline-color,-moz-outline-style,
    -moz-outline-width,-moz-user-focus,-moz-user-input,
    -moz-user-modify,-moz-user-select,orphans,outline,
    outline-color,outline-style,outline-width,overflow,
    overflow-X,overflow-Y,padding,padding-bottom,padding-left,
    padding-right,padding-top,page,page-break-after,
    page-break-before,page-break-inside,pause,pause-after,
    pause-before,pitch,pitch-range,play-during,position,quotes,
    -replace,richness,right,ruby-align,ruby-overhang,
    ruby-position,-set-link-source,size,speak,speak-header,
    speak-numeral,speak-punctuation,speech-rate,stress,
    scrollbar-arrow-color,scrollbar-base-color,
    scrollbar-dark-shadow-color,scrollbar-face-color,
    scrollbar-highlight-color,scrollbar-shadow-color,
    scrollbar-3d-light-color,scrollbar-track-color,table-layout,
    text-align,text-align-last,text-decoration,text-indent,
    text-justify,text-overflow,text-shadow,text-transform,
    text-autospace,text-kashida-space,text-underline-position,top,
    unicode-bidi,-use-link-source,vertical-align,visibility,
    voice-family,volume,white-space,widows,width,word-break,
    word-spacing,word-wrap,writing-mode,z-index,zoom},
  morestring=[s]{:}{;},
  sensitive,
  morecomment=[s]{/*}{*/}
}

% Spolszczenia nazw elementów
\renewcommand{\lstlistingname}{Listing}
\renewcommand{\lstlistlistingname}{Spis listingów}
\renewcommand{\listfigurename}{Spis rysunków}
\renewcommand{\listtablename}{Spis tabel}
\renewcommand{\figurename}{Rysunek}
\renewcommand{\tablename}{Tabela}
\renewcommand{\chaptername}{Rozdział}
\renewcommand{\contentsname}{Spis Treści}

% Środowisko do warunków (np. w matematyce)
\newenvironment{conditions}
  {\par\vspace{\abovedisplayskip}\noindent\begin{tabular}{>{$}l<{$} @{${}={}$} l}}
  {\end{tabular}\par\vspace{\belowdisplayskip}}



% --- NOWE ŚRODOWISKO "GDZIE" (opis do wzorow) ---
\newenvironment{gdzie}
{
    \vspace{-1.0em}    % Podciągnij tekst do góry (bliżej wzoru)
    \noindent \textbf{gdzie:}  % Automatyczny napis "gdzie:"
    \begin{itemize}[         % Lista z ciasnymi ustawieniami
        noitemsep,           % Brak odstępów między punktami
        topsep=2pt,  % Mały odstęp od "gdzie:" do pierwszego punktu
        parsep=0pt,          % Brak odstępów wewnątrz punktów
        partopsep=0pt,       % Brak dodatkowych odstępów
        leftmargin=1.5em,    % Wcięcie listy
        labelsep=0.5em     % Odstęp kropki/myślnika od tekstu
    ]
}
{
    \end{itemize}
    \vspace{0.5em}                  % Odstęp po zakończeniu opisu
}